\documentclass[11pt]{article}

\usepackage[margin=25mm]{geometry}

\usepackage{amsmath,amssymb,fontspec,unicode-math}

\IfFontExistsTF{DejaVu Serif}{\setmainfont{DejaVu Serif}}{\setmainfont{Latin Modern Roman}}

\IfFontExistsTF{Latin Modern Math}{\setmathfont{Latin Modern Math}}{\setmathfont{DejaVu Math TeX Gyre}}

\usepackage{xurl,hyperref}

\setlength{\parindent}{0pt}\setlength{\parskip}{6pt}

\begin{document}

\title{Semantic Representation Attack against Aligned Large Language Models}

\author{ELVIA preliminary AI analysis}\date{21 October 2026}\maketitle

Preliminary assessment: Major.\par

Independent scientific review remains pending. Author adoption has not been requested.

The reciprocal probability theorem can demand a probability greater than one. Representative/class optimization, infinite-class floors and nucleus sampling need narrower conclusions, while the original fixed-floor perplexity sufficient condition survives.

\section*{Review scope}

Latest-version, claim-focused review of the cited author-paper passages and applicable current author-release paths. The stated counterexamples and scope controls were checked analytically; no author program, old verifier, training run or benchmark was executed.

\section*{Problem 1: The reciprocal lower bound can exceed one}

Source: Theorem 3.1, Eq.(5), PDF pp.4–5; Appendix A.1, PDF pp.24–25, Eqs.(19)–(31)

Assessment: supported; impact: major

\subsection*{Source claim (paraphrase)}

Theorem 3.1 claims \(P(\hat y_2\mid c)>\delta/(D_{\rm seq}+\epsilon)\) under its coherence and semantic-class assumptions.

\subsection*{Analysis}

Let one-token context \(c\) have probability \(1\), followed by completed responses \(a,b,r\) with probabilities \((0.4,0.4,0.2)\), then certain termination. Put \(a,b\) in the same semantic class, \(\tau=2\), \(\epsilon=0.01\). Coherence is \(\sqrt{2.5}<2\); the appendix gives \(\delta=1/4\), while \(D_{\rm seq}=\log(0.4/0.4)=0\). The theorem demands \(0.4>25\), despite meeting its small-argument condition. Linearizing an exponential and then dropping the added constant does not preserve a numerical lower bound.

\subsection*{Finding and scope}

The central reciprocal theorem is false under its own assumptions. No model or operational attack is needed.

\subsection*{Proposed correction}

Retain the exact log-ratio identity \(P(\hat y_2\mid c)=P(\hat y_1\mid c)e^{-D_{\rm seq}}\), which gives lower bound \(\delta e^{-D_{\rm seq}}\). It does not prove semantic convergence by itself.

An explicit local correction is described; author adoption and a corresponding published erratum are not established.

\section*{Problem 2: An improved representative can reduce semantic-class mass}

Source: Theorem 3.2, Eq.(6), PDF p.6; Appendix A.2, PDF pp.26–27

Assessment: supported; impact: minor

\subsection*{Source claim (paraphrase)}

Theorem 3.2’s argument treats maximizing one representative as maximizing the class mass, without a stated approximation factor or fixed within-class proportions.

\subsection*{Analysis}

Two feasible contexts give class probabilities \((0.400,0.380)\) and \((0.395,0.395)\) for responses \(a,b\). Maximizing \(a\) selects the first, but class masses are \(0.780<0.790\). Embed the contexts with probabilities \(0.005\) each, certain response termination, \(\tau=100\) and \(\epsilon=0.01\); then \(\delta=0.0002\), all class probabilities are coherent and all ordered-pair small-ratio conditions hold. The monotonic implication fails independently of the first theorem.

\subsection*{Finding and scope}

Optimizer equivalence and the proof’s monotonicity are unsupported. An unspecified empirical approximation tolerance is not quantitatively refuted.

\subsection*{Proposed correction}

Fixed positive proportions \(P(y_i\mid c)=w_iP(Y_\Phi\mid c)\), with \(w_i\) independent of context, are a sufficient condition for exact equivalence. They are not established properties of the models.

An explicit local correction is described; author adoption and a corresponding published erratum are not established.

\section*{Problem 3: An infinite completed-response class cannot share a positive floor}

Source: Definition 3.1, Eq.(2), PDF p.3; sequence domain in Section 3.1

Assessment: supported; impact: minor

\subsection*{Source claim (paraphrase)}

Definition 3.1 requires \(P(y_i\mid c)>\delta>0\) for every member of a potentially infinite semantic class.

\subsection*{Analysis}

Distinct completed responses are disjoint events. For any \(N\) members, normalization gives \(1\geq\sum_iP(y_i\mid c)>N\delta\), impossible once \(N>1/\delta\). A finite class requires \(N\delta<1\). Overlapping prefixes cannot instead be summed as completed-response mass.

\subsection*{Finding and scope}

The formulation needs a finite feasible domain or a different coverage objective. The implemented finite search is not shown empty.

\subsection*{Proposed correction}

Specify a finite candidate class with compatible thresholds, or use a class-mass constraint. These are different objectives and do not automatically preserve Theorem 3.2.

An explicit local correction is described; author adoption and a corresponding published erratum are not established.

\section*{Problem 4: Nucleus mass does not enforce a fixed token-probability floor}

Source: Theorem 3.3 and Algorithm 1; Appendix B.2, PDF pp.28–29; baselines/sra/utils.py, GenerationWorker selector

Assessment: supported; impact: minor

\subsection*{Source claim (paraphrase)}

Theorem 3.3 uses per-token probability above \(1/\tau\), while the released selector uses cumulative top-p mass and top-k truncation.

\subsection*{Analysis}

Probabilities \((0.90,0.04,0.03,0.03)\) at top-p level \(0.95\) select the first three tokens. At \(\tau=20\), two are below \(0.05\). A one-token context of perplexity \(19.9\) followed by the \(0.03\) token has perplexity \(\sqrt{19.9/0.03}>20\). The selector contains no fixed-floor check.

\subsection*{Finding and scope}

Nucleus selection alone does not preserve the same fixed-threshold guarantee. The original sufficient theorem remains valid.

\subsection*{Proposed correction}

Enforce the fixed floor or check cumulative perplexity directly. An adaptive threshold needs an explicit update and final validation before claiming its changed guarantee.

An explicit local correction is described; author adoption and a corresponding published erratum are not established.

\section*{Limits}

Preliminary AI assessment; independent human scientific adjudication remains pending.\par

An analytic counterexample or a release discrepancy does not by itself remeasure the paper’s reported experiments.\par

\section*{Source identity}

NeurIPS 2025 official proceedings PDF; compared with current arXiv:2509.19360v2

\url{https://proceedings.neurips.cc/paper_files/paper/2025/file/6958e9d0f5a76d54ff97da8c45f4d52e-Paper-Conference.pdf}

Reviewed PDF SHA-256: \nolinkurl{70722ef94b2886869a9e71c0cf7a145e162b7351bd6e452b0ea58e37ed4fade7}

The selected author/publisher source is the same work as the original case, matched by exact title, author identity and official version channel. The exact selected PDF SHA-256 was recomputed in the latest-version pass for Issue \#27. The PDF is a retained author/publisher copy; its original retrieval time is stated only when supported by matching acquisition metadata. File modification time was not substituted.

\end{document}
